% =====================================================================
% 5  Heterogeneity                     Folie 14 | 1:45 min | bis 18:00
% Inhalt: Bericht Kap. 8. Pflicht laut instruction.pdf: WLS-Meta-
% Regression zu Publication bias und (falls moeglich) Heterogenitaet.
% =====================================================================
\section{Sources of heterogeneity}

% ---------------------------------------------------------------------
\begin{frame}{Sources of heterogeneity: WLS meta-regression}
  \timing{1:45}{18:00}
  \vspace{-0.6em}
  {\small\[
    y_i = \beta_0 + \beta_1 SE_i + \textstyle\sum_k \delta_k Z_{ik}
      + \varepsilon_i, \qquad w_i = 1/SE_i^2
  \]}
  \vspace{-1.2em}
  \slidetablebare{metareg_wls}{0.6\textheight}
  \keymsg{Europe: stronger discrimination. Controlling for region shrinks the
    FAT (M1 $\to$ M2) $\Rightarrow$ part of the asymmetry is heterogeneity.}
\end{frame}
% Sprechnotizen (Folie 14) -- Zahlen aus der Tabelle, nicht auf der Folie getippt:
% - Lesen: Koeffizienten auf der ln-OR-Skala, negativ = mehr Diskriminierung.
%   Referenzgruppe: ausserhalb Europas, gemischte Skills, Testjahr 2000.
% - Europe: -0.28 (SE 0.13) in M2 -> OR ca. 25 % tiefer; stabil in M2-M5.
% - FAT faellt von -3.16 (M1) auf -1.91 (M2): europaeische Studien sind im
%   Schnitt unpraeziser UND finden mehr Diskriminierung -> ein Teil der
%   Funnel-Asymmetrie ist Zusammensetzung der Stichprobe.
% - Nicht signifikant: German-speaking (+0.20), Skills (M3, nur k = 30),
%   Zeittrend (-0.01), Dummy neue Studien (0.04).
% - Caveat: nur 38 Studien, korrelierte Moderatoren, vieles bleibt unerklaert.
